\section{Where a result becomes public after interruption}
\label{exec17:i4-publication}

A recovery specification should identify the transition at which a consumer
may rely on a result. The calculation finishing, the artifact bytes appearing,
and the job becoming completed are distinct events. In the inspected local
report worker, a job lock protects the final cancellation check, artifact
publication and update of the job's completed state and result reference.
The artifact store has its own content and metadata publications. This
source structure gives a concrete basis for reasoning about selected crash
points; it does not create a transaction over every file or an external
service \citep{exec13_ops_source}.

Consider the sequence $C\to A\to T$, where $C$ is a validated completed-work
checkpoint, $A$ means the report artifact has been stored, and $T$ means the
job record exposes a completed state and the artifact reference. Before $C$,
uncommitted work may need recalculation. Between $C$ and $A$, recovery can
reconstruct a report from the checkpoint after revalidating identities.
Between $A$ and $T$, an artifact may exist without being exposed by a
completed job. After $T$, a client that loses the response should retrieve
the result associated with that identity rather than create a new analysis.
These states explain why counting worker starts does not count distinct
scientific results, and why the existence of a report file alone should not
be the consumer's completion criterion.

The actual report path uses deterministic report contents and a content-based
artifact identity. A repeated publication checks the same content rather
than appending a second result to the public report. That supports recovery
for this path. A future estimator with stochastic or environment-sensitive
output would need a preserved randomization and numerical identity before
the same argument could apply. A new draw after a crash may be a new
statistical replicate; it is not automatically a continuation of the original
one. Likewise, byte-identical output can still embody false causal premises.
Recovery identity and scientific validity answer different questions.

Cancellation can be analyzed by the order of the protected transitions. In a
small authored state machine, begin at running. If cancellation is accepted
first, the state becomes cancelled and publication is disallowed. If the
publication transition completes first, the state becomes completed with a
public result and a later cancellation leaves that state intact. Enumeration
of these two orders reproduces the two outcomes
\citep{exec17_integration_review_reference}. The example assumes an indivisible
protected transition and a shared authoritative state. It is a reference
model for the local lock boundary, not a proof for independent machines,
uncooperative numerical kernels or external banking actions.

The lock also clarifies what a consumer can reasonably be told. A cancellation
request can wait while publication holds the lock. The returned completed
state then reports what actually happened; it must not claim that an earlier
user gesture erased the report. Conversely, a process killed while holding
the lock can leave an artifact published before the terminal state is
written. Recovery must recheck the job and artifact identities rather than
infer completion from the interrupted caller's intention. The preserved
fault fixtures exercise selected local interruptions, but they do not
enumerate every instruction boundary or prove crash consistency for all
filesystems and hardware \citep{exec13_ops_faults}.

A complete prototype fault protocol should state its injection points and
the invariants inspected after each restart. For the report path those
include: the same request and source identity; a valid completed-work prefix;
one contribution for each committed work identity; a valid artifact reference
only in a completed public state; and a truthful terminal cancellation
outcome. A negative case should change an input, checkpoint or environment
identity and require a refusal or new analysis rather than continuation.
These are proposal acceptance conditions. A platform study must additionally
state what durability its storage actually guarantees when the whole host,
not merely one worker, fails.

\subsection{Idempotence does not cross an uncooperative boundary}

An external offer illustrates the limit. Suppose an adapter sends a customer
message and loses the acknowledgement before recording success. There are
two worlds consistent with the local observation: the message was delivered,
or it was not. Blind retry duplicates the first world's action; refusing to
retry omits the second world's action. A local content hash cannot distinguish
them. A cooperating recipient that recognizes a durable action identity,
or an authoritative query that resolves delivery, can change the information
available. Without such a contract, the correct state is unresolved and a
reconciliation owner is needed. The current report service performs no such
action, so this example motivates a later interface requirement rather than
alleging an observed product failure.

This distinction should affect product scope and cost. A read-oriented report
prototype can demonstrate semantic preservation and local recovery without
having permission to act on customer accounts. An action-taking extension
adds authorization, recipient deduplication, reconciliation and compensation
requirements that must be costed and tested independently. A compensating
action, when one exists, is a new action with its own consequences; it is
not evidence that the original event never occurred. The proposal should
retain these boundaries when comparing workflow automation alternatives.

\section{Retention is reachability under a stable ownership decision}
\label{exec17:i4-retention}

Immutable artifacts make recovery auditable but create a storage obligation.
Deleting every file older than a chosen age can remove evidence still needed
by a current report or a resumed job. Keeping every file forever avoids that
particular mistake but makes cost and capacity unbounded in an ongoing
service. A proposed retention design should therefore distinguish the
scientific dependency graph from the policy that designates live roots.
The former records which exact objects are needed; the latter states which
reports, jobs, reviews and retention holds must remain available.

Let $\mathcal A$ be the finite set of stored artifact identities and let
$a\to b$ mean that artifact $a$ depends on $b$. Given a root set
$\mathcal R$, define the retained closure recursively by
\begin{equation}
 K_0=\mathcal R,\qquad
 K_{j+1}=K_j\cup\{b:\exists a\in K_j,\ a\to b\},\qquad
 K=\bigcup_{j\geq0}K_j.
 \label{exec17:i4-reachability}
\end{equation}
In a finite graph this process stabilizes after at most $|\mathcal A|$
strict additions of new identities, and cycles do not prevent termination
when visited identities are recorded. The deletion candidates are
$\mathcal A\setminus K$. This is a graph statement under complete dependency
metadata. It does not decide legal retention, scientific adequacy, access
authorization, or whether an owner has supplied every required root.

The exact reference has live report and checkpoint roots. Both reach a result
that depends on data and code. The retained set therefore contains those
five identities. An orphan artifact and an old result with its old input
form three deletion candidates when no other root references them.
The reference traverses the graph and checks this result, without deleting
any files \citep{exec17_integration_review_reference}. Adding an archived
review as a root could legitimately retain the old result. Removing an
unfinished job from the roots merely because its process is absent could
destroy recovery evidence. Root ownership is consequently a substantive
operating decision, not an incidental garbage-collector parameter.

Concurrent publication complicates the simple closure. A collector can
observe an artifact before a new job record points to it, classify it as
unreachable, and delete it just as the publisher exposes the reference.
A safe proposed implementation needs a stable root snapshot and coordination
with publication, or another protocol that prevents such a race. A
quarantine interval alone has no proof without an upper bound on publication
and reconciliation delays. The prototype can compare a lock-based local
collection phase with a versioned snapshot and explicit publication pins.
Each alternative should state the failure model and verify that a reachable
artifact is never deleted under that model.

The dependency graph should also preserve descriptor identities and custody.
A report referencing a small summary may still depend on a raw input that
must remain in a restricted store. Reachability does not require making all
reachable content readable to the same client. The proposed collector
operates under the store owner's authority, while a reader still passes
role, size and custody checks. Similarly, duplicate content can be stored
once physically while retaining distinct provenance or policy descriptors;
equating physical deduplication with interchangeable scientific evidence
would discard information about use and interpretation.

For a transparent storage calculation, let $n$ be admitted requests per day,
$s$ the mean additional retained bytes per request under a specified policy,
and $h$ the number of retained days. Ignoring shared content, metadata,
temporary files and replication gives $nsh$ bytes. This is an accounting
baseline, not a capacity guarantee. A real measurement must add those
categories and account for distributions: a few large covariance artifacts
can dominate a mean-sized planning example. Retention holds and historical
reviews can also extend effective lifetime beyond $h$. The proposal should
show which quantities are measured locally, stipulated for planning, or
awaiting a target operator's requirements.

\section{A worker limit is not an admission or deadline guarantee}
\label{exec17:i4-capacity}

The inspected worker has a memory and wall-time limit, while its coordinator
and SDK consume additional resources. An operator therefore needs a
whole-workflow admission rule. Suppose a hypothetical host budget is
1024 MiB, the resident parent requires 128 MiB, a reserve takes another
128 MiB, and each admitted worker is allotted 256 MiB. Under an additive
worst-case model the maximum concurrent count is
\begin{equation}
 k=\left\lfloor\frac{1024-128-128}{256}\right\rfloor=3.
 \label{exec17:i4-memory}
\end{equation}
Four workers would fit only if the parent or reserve were ignored. These
figures are deliberately stipulated and reproduced by the reference
\citep{exec17_integration_review_reference}. Actual address-space limits are
not measurements of peak resident memory, and shared pages, buffers,
serialization and numerical libraries can alter the relationship. The
equation describes a conservative admission model once its quantities and
scope have been justified; it does not qualify the present host.

Admission also affects delay before work starts. Consider five simultaneous
requests, two identical workers, and exactly three seconds of service per
request. With a first-available-worker rule, completion times are
$3,3,6,6,9$ seconds. A deadline eight seconds after arrival is met by four
requests and missed by one, even though every service lasts only three
seconds. Direct scheduling and the independent expression
$3\lceil j/2\rceil$, $j=1,\ldots,5$, agree in the exact reference.
Thus a per-worker time limit can hold while the user-facing deadline fails.
The denominator includes the fifth request. Dropping it as an unfinished
job would overstate deadline performance.

The relevant deadline must be declared at intake. If $t_i$ is arrival time
and $b_i$ is the allowed end-to-end duration, a later start at $s_i$ leaves
at most $t_i+b_i-s_i$ for remaining work, less any required publication and
delivery allowance. Starting a fresh full budget after queueing changes the
contract. If the remaining allowance is already exhausted, rejecting or
expiring the request is more truthful than launching work and relabeling
its eventual result on-time. If the service-time bound is unknown, a queue
policy can report an estimate with its limitations; it cannot provide a
hard completion promise merely by limiting concurrency.

The queue example intentionally removes random arrivals, heterogeneous
service times and failures. Those extensions need a stated workload law or
measured trace, not a casual independence assumption. Correlated daily
campaigns can create bursts even when the long-run average arrival rate
looks modest. Long jobs can delay short ones under one policy; prioritizing
short jobs can systematically exclude the scientifically difficult tasks.
The proposal should compare admission and scheduling rules using complete
task identities and the same mix of targets, with rejected, expired and
cancelled requests visible. A faster successful subset is not a generally
faster workflow.

Back-pressure is therefore a scientific-interface issue as well as an
operating concern. A rejected request needs a distinct execution state and
its original target. It must not become an unsupported causal-method claim.
A timeout during a difficult latent-graph calculation cannot be counted as
proof of nonidentification. Conversely, an identified but unbounded result
should not be retried until it becomes a convenient scalar. Capacity,
scientific disposition and usefulness should be recorded in separate fields
so that each failure has an appropriate remedy.

\subsection{What an inexpensive complete trial should measure}

A public/generated prototype trial can now be specified without pretending
to have production data. Freeze a finite mixture of supported and negative
scientific requests, artifact sizes and budget classes. Record intake,
admission, queue entry, worker start, computation completion, publication
and client receipt under one clock convention. Measure peak resources for
the process tree and store, include failed parsing and serialization, and
inject declared crashes before and after checkpoint and publication. The
retention trial should operate on a disposable copy with known dependencies
and verify every live reference after collection. The reference constructions
in this section make expected states auditable; they do not substitute for
the subsequent measurements.

Compare the library, CLI and compatible MCP path under the same request
definitions, while separately reporting their unavoidable launch and
transport costs. Use an incumbent analyst or scripted workflow as a
decision comparator only when its task and permitted information match.
Monetary costs should distinguish one-time construction, shared operation,
per-request compute, review, retry and recovery. Avoid counting a worker's
CPU charge again inside an already inclusive service charge. The prior
all-attempt example shows why costs from failures and fallback decisions
belong in the economic denominator \citep{exec13_ops_exact}.

The resulting evidence can support a scoped design decision: which consumer
boundary preserves meaning, which storage rule maintains recovery, and what
workload the local prototype can actually process. It cannot authenticate a
bank owner, supply real assignment truth or establish institution-scale
availability. Those later trials remain conditional on prototype merit and
agreed acceptance criteria. Maintaining that distinction permits the present
proposal to finish its engineering argument while keeping the expensive
external commitments explicit.
\label{exec17:i4-capacity-end}
